\documentclass[12pt,a4paper]{report}

\usepackage{fontspec}
\usepackage{polyglossia}
\setmainlanguage{romanian}
\setmainfont{Latin Modern Roman}

\usepackage[a4paper,margin=2.5cm]{geometry}
\usepackage{setspace}
\usepackage{amsmath}
\usepackage{amssymb}
\usepackage{enumitem}
\usepackage{longtable}
\usepackage{array}
\usepackage{booktabs}
\usepackage{xcolor}
\usepackage{listings}
\usepackage[hidelinks]{hyperref}

\onehalfspacing
\setlength{\parindent}{1.2cm}
\setlength{\parskip}{0.25em}

\definecolor{codebg}{RGB}{248,248,248}
\definecolor{codeframe}{RGB}{210,210,210}
\definecolor{codekeyword}{RGB}{0,78,130}
\definecolor{codestring}{RGB}{120,70,0}
\definecolor{codecomment}{RGB}{80,120,80}

\lstdefinestyle{pythonstyle}{
    language=Python,
    basicstyle=\ttfamily\small,
    keywordstyle=\color{codekeyword}\bfseries,
    stringstyle=\color{codestring},
    commentstyle=\color{codecomment}\itshape,
    backgroundcolor=\color{codebg},
    frame=single,
    rulecolor=\color{codeframe},
    numbers=left,
    numberstyle=\tiny,
    numbersep=8pt,
    breaklines=true,
    showstringspaces=false,
    tabsize=4,
    keepspaces=true
}

\newcommand{\codechunk}[2]{
    \lstinputlisting[
        style=pythonstyle,
        firstline=#1,
        lastline=#2,
        firstnumber=#1
    ]{../src/dsp/metrics.py}
}

\title{Explicatie detaliata a fisierului \texttt{src/dsp/metrics.py}}
\author{Documentatie tehnica pentru proiectul de procesare audio}
\date{\today}

\begin{document}

\maketitle
\tableofcontents

\chapter{Rolul fisierului in proiect}

Fisierul \texttt{src/dsp/metrics.py} contine functiile prin care proiectul masoara calitatea rezultatelor audio si performanta detectorului VAD. Daca fisierele precum \texttt{spectral\_subtraction.py}, \texttt{wiener\_filter.py} si \texttt{inference.py} produc semnale imbunatatite, acest fisier raspunde la intrebarea: cat de buna este imbunatatirea?

Metricile implementate acopera doua zone:

\begin{itemize}
    \item evaluarea speech enhancement: SNR, STOI, PESQ, SI-SDR, segmental SNR si log spectral distance;
    \item evaluarea VAD: accuracy, precision, recall, F1, miss rate si false alarm rate.
\end{itemize}

Acest modul este folosit ca infrastructura de evaluare. El permite compararea intre semnalul zgomotos, metodele clasice DSP si MaskNet.

\chapter{Importuri si dependinte optionale}

\section{Liniile 1--14}

\codechunk{1}{14}

\begin{description}[leftmargin=2.9cm,style=nextline]
    \item[Linia 1] Activeaza evaluarea amanata a adnotarilor de tip, pentru compatibilitate moderna.
    \item[Linia 2] Importa \texttt{Dict}, folosit pentru tipul dictionarelor de metrici.
    \item[Linia 3] Importa \texttt{numpy}, folosit pentru calcule numerice.
    \item[Linia 4] Importa \texttt{config}, de unde se ia rata de esantionare implicita.
    \item[Liniile 6--9] Incearca sa importe functia STOI din pachetul \texttt{pystoi}. Daca pachetul nu este instalat, variabila \texttt{\_stoi} devine \texttt{None}.
    \item[Liniile 11--14] Incearca sa importe functia PESQ din pachetul \texttt{pesq}. Daca nu exista, \texttt{\_pesq} devine \texttt{None}.
\end{description}

Dependintele STOI si PESQ sunt tratate optional. Modulul poate fi incarcat chiar daca aceste pachete nu sunt instalate, iar eroarea apare doar cand se incearca folosirea metricii respective.

\chapter{SNR global}

\section{Liniile 16--36}

\codechunk{16}{36}

Functia \texttt{snr\_db} calculeaza raportul semnal-zgomot in decibeli:

\[
SNR_{dB} = 10\log_{10}\frac{P_{semnal}}{P_{zgomot}},
\]

unde:

\[
P_{semnal} = \frac{1}{N}\sum_n x[n]^2,
\quad
P_{zgomot} = \frac{1}{N}\sum_n (x[n]-\hat{x}[n])^2.
\]

\begin{description}[leftmargin=2.9cm,style=nextline]
    \item[Linia 16] Defineste functia pentru semnalul curat si semnalul procesat.
    \item[Liniile 17--19] Verifica daca vreun vector este gol. In acest caz ridica \texttt{ValueError}.
    \item[Liniile 21--24] Daca semnalele au mai multe canale, le converteste la mono prin medie.
    \item[Linia 25] Alege lungimea comuna minima.
    \item[Liniile 27--28] Verifica daca exista suprapunere intre semnale.
    \item[Liniile 30--31] Taie semnalele la aceeasi lungime si le converteste la \texttt{float32}.
    \item[Linia 32] Calculeaza eroarea fata de referinta, interpretata ca zgomot/distorsiune.
    \item[Liniile 33--35] Calculeaza puterea semnalului si puterea erorii, folosind \(\varepsilon\) pentru stabilitate.
    \item[Linia 36] Returneaza SNR-ul in dB.
\end{description}

In acest proiect, un SNR mai mare inseamna ca semnalul procesat este mai apropiat energetic de semnalul curat.

\chapter{STOI si PESQ}

\section{STOI: liniile 38--46}

\codechunk{38}{46}

STOI inseamna \emph{Short-Time Objective Intelligibility}. Este o metrica folosita pentru inteligibilitatea vorbirii.

\begin{description}[leftmargin=2.9cm,style=nextline]
    \item[Linia 38] Defineste functia \texttt{stoi\_score}.
    \item[Liniile 39--40] Daca \texttt{pystoi} nu este instalat, ridica \texttt{ImportError}.
    \item[Liniile 41--42] Daca rata de esantionare nu este transmisa, foloseste \texttt{config.SAMPLE\_RATE}.
    \item[Liniile 43--45] Aliniaza semnalele la aceeasi lungime si le converteste la \texttt{float32}.
    \item[Linia 46] Apeleaza implementarea STOI si returneaza scorul ca \texttt{float}.
\end{description}

\section{PESQ: liniile 48--56}

\codechunk{48}{56}

PESQ inseamna \emph{Perceptual Evaluation of Speech Quality}. Este o metrica perceptuala pentru calitatea vorbirii.

\begin{description}[leftmargin=2.9cm,style=nextline]
    \item[Linia 48] Defineste functia \texttt{pesq\_score}.
    \item[Liniile 49--50] Daca pachetul \texttt{pesq} nu este instalat, ridica \texttt{ImportError}.
    \item[Liniile 51--52] Stabileste rata de esantionare implicita.
    \item[Liniile 53--55] Aliniaza semnalele si le converteste la \texttt{float32}.
    \item[Linia 56] Apeleaza PESQ in modul wideband, marcat prin argumentul \texttt{"wb"}.
\end{description}

STOI si PESQ sunt importante pentru ca merg dincolo de eroarea matematica simpla si incearca sa aproximeze inteligibilitatea si calitatea perceptuala.

\chapter{Functii wrapper si SI-SDR}

\section{Wrapper-e de compatibilitate: liniile 58--69}

\codechunk{58}{69}

\begin{description}[leftmargin=2.9cm,style=nextline]
    \item[Linia 58] Comentariul explica rolul functiilor: compatibilitate cu alte parti ale proiectului.
    \item[Liniile 59--61] \texttt{compute\_snr} apeleaza \texttt{snr\_db}.
    \item[Liniile 63--65] \texttt{compute\_stoi} apeleaza \texttt{stoi\_score}.
    \item[Liniile 67--69] \texttt{compute\_pesq} apeleaza \texttt{pesq\_score}.
\end{description}

Aceste functii permit altor fisiere sa foloseasca denumiri standardizate fara sa depinda de numele interne.

\section{SI-SDR: liniile 71--90}

\codechunk{71}{90}

SI-SDR inseamna \emph{Scale-Invariant Signal-to-Distortion Ratio}. Ideea este sa se proiecteze semnalul procesat pe semnalul curat, astfel incat diferenta de scala sa nu domine scorul.

\[
\alpha = \frac{\hat{x}^T x}{x^T x + \varepsilon},
\quad
s_{target} = \alpha x,
\quad
e = \hat{x} - s_{target}.
\]

Scorul este:

\[
SI\text{-}SDR =
10\log_{10}
\frac{\sum_n s_{target}[n]^2}
{\sum_n e[n]^2+\varepsilon}.
\]

\begin{description}[leftmargin=2.9cm,style=nextline]
    \item[Liniile 71--72] Definesc functia si docstring-ul.
    \item[Liniile 73--76] Convertesc semnalele multi-canal la mono.
    \item[Liniile 78--80] Aliniaza semnalele la aceeasi lungime.
    \item[Linia 83] Calculeaza factorul de proiectie \(\alpha\).
    \item[Linia 84] Calculeaza componenta tinta scalata.
    \item[Linia 85] Calculeaza eroarea/distorsiunea ramasa.
    \item[Liniile 88--90] Calculeaza si returneaza scorul SI-SDR.
\end{description}

\chapter{Evaluarea VAD}

\section{Liniile 92--116}

\codechunk{92}{116}

Functia \texttt{evaluate\_vad} compara etichetele reale VAD cu etichetele prezise.

\begin{description}[leftmargin=2.9cm,style=nextline]
    \item[Liniile 92--94] Defineste functia si converteste etichetele la intregi.
    \item[Liniile 95--97] Aliniaza vectorii la aceeasi lungime.
    \item[Liniile 98--101] Calculeaza valorile matricei de confuzie: true positive, true negative, false positive si false negative.
    \item[Linia 102] Defineste o constanta mica pentru stabilitate.
    \item[Linia 103] Calculeaza acuratetea.
    \item[Linia 104] Calculeaza precizia.
    \item[Linia 105] Calculeaza recall-ul.
    \item[Linia 106] Calculeaza scorul F1.
    \item[Linia 107] Calculeaza rata de ratari, adica proportia cadrelor vocale neidentificate.
    \item[Linia 108] Calculeaza rata alarmelor false, adica proportia cadrelor de liniste marcate ca vorbire.
    \item[Liniile 109--116] Returneaza toate metricile intr-un dictionar.
\end{description}

Formulele principale sunt:

\[
Precision = \frac{TP}{TP+FP},
\quad
Recall = \frac{TP}{TP+FN},
\quad
F1 = \frac{2PR}{P+R}.
\]

\chapter{Segmental SNR}

\section{Liniile 119--152}

\codechunk{119}{152}

\texttt{segmental\_snr} calculeaza SNR-ul pe cadre scurte si apoi mediaza valorile. Aceasta metoda evita ca un singur segment foarte energetic sa domine scorul global.

\begin{description}[leftmargin=2.9cm,style=nextline]
    \item[Liniile 119--120] Definesc functia si docstring-ul.
    \item[Liniile 121--128] Convertesc la mono, aliniaza lungimile si convertesc semnalele la \texttt{float32}.
    \item[Linia 130] Calculeaza eroarea dintre semnalul curat si cel procesat.
    \item[Linia 131] Calculeaza numarul de cadre complete.
    \item[Liniile 133--134] Daca nu exista niciun cadru complet, returneaza \texttt{0.0}.
    \item[Liniile 136--150] Parcurg cadrele, calculeaza SNR-ul fiecaruia si il limiteaza intre \(-10\) si \(35\) dB.
    \item[Linia 152] Returneaza media valorilor pe cadre.
\end{description}

Limitarea valorilor pe cadre este o practica utila deoarece cadrele aproape tacute pot produce valori extreme.

\chapter{Log Spectral Distance}

\section{Liniile 155--180}

\codechunk{155}{180}

Log Spectral Distance masoara distanta dintre spectrele logaritmice ale semnalului curat si semnalului procesat:

\[
LSD_t =
\sqrt{
\frac{1}{F}
\sum_f
\left(
10\log_{10}P_x(f,t)
-
10\log_{10}P_{\hat{x}}(f,t)
\right)^2
}.
\]

\begin{description}[leftmargin=2.9cm,style=nextline]
    \item[Liniile 155--156] Definesc functia si docstring-ul.
    \item[Liniile 157--164] Convertesc la mono, aliniaza lungimile si convertesc la \texttt{float32}.
    \item[Linia 167] Importa local \texttt{scipy.signal}. Importul local evita dependenta directa pana cand functia este folosita.
    \item[Linia 168] Stabileste lungimea segmentului spectrogramelor.
    \item[Liniile 169--170] Calculeaza spectrogramele de putere pentru semnalul curat si procesat.
    \item[Liniile 173--175] Aliniaza dimensiunea temporala a celor doua spectrograme.
    \item[Liniile 178--179] Calculeaza LSD pe fiecare cadru.
    \item[Linia 180] Returneaza media distantei spectrale.
\end{description}

Pentru LSD, valori mai mici inseamna spectru mai apropiat de referinta curata.

\chapter{Evaluatorul agregat}

\section{Liniile 183--230}

\codechunk{183}{230}

Functia \texttt{evaluate\_speech\_enhancement} este interfata principala pentru evaluarea unui semnal imbunatatit fata de referinta curata.

\begin{description}[leftmargin=2.9cm,style=nextline]
    \item[Liniile 183--187] Definesc functia pentru semnalul curat, semnalul imbunatatit si rata de esantionare.
    \item[Liniile 188--198] Documenteaza scopul: evaluare completa cu PESQ, STOI, SNR, SegSNR si LSD.
    \item[Liniile 199--200] Stabileste rata de esantionare implicita.
    \item[Linia 202] Creeaza dictionarul de rezultate.
    \item[Liniile 205--208] Incearca sa calculeze PESQ. Daca apare o eroare, pune \texttt{0.0}.
    \item[Liniile 210--213] Calculeaza STOI cu aceeasi strategie de fallback.
    \item[Liniile 215--218] Calculeaza SNR.
    \item[Liniile 220--223] Calculeaza segmental SNR.
    \item[Liniile 225--228] Calculeaza LSD.
    \item[Linia 230] Returneaza dictionarul de metrici.
\end{description}

Strategia cu \texttt{try/except} face ca evaluarea sa nu se opreasca complet daca o metrica nu este disponibila sau esueaza pentru un fisier. Dezavantajul este ca valoarea \texttt{0.0} poate ascunde motivul exact al erorii, deci pentru analiza detaliata trebuie consultate si logurile din scripturile de evaluare.

\chapter{Fluxul complet al fisierului}

Fluxul conceptual este:

\begin{enumerate}
    \item semnalul curat si semnalul procesat sunt aduse la aceeasi lungime;
    \item daca sunt stereo, sunt convertite la mono;
    \item se calculeaza metrici energetice: SNR, segmental SNR, SI-SDR;
    \item se calculeaza metrici perceptuale, daca dependintele sunt instalate: STOI si PESQ;
    \item se calculeaza distanta spectrala logaritmica;
    \item pentru VAD, se calculeaza matricea de confuzie si metricile derivate.
\end{enumerate}

\chapter{Concluzie}

\texttt{src/dsp/metrics.py} este fisierul care transforma rezultatele audio in valori masurabile. El este esential pentru compararea metodelor, pentru raportarea performantei si pentru validarea faptului ca MaskNet sau metodele DSP chiar imbunatatesc semnalul.

Prin includerea metricilor SNR, STOI, PESQ, SI-SDR, segmental SNR, LSD si VAD, modulul acopera atat apropierea matematica fata de semnalul curat, cat si aspecte legate de inteligibilitate, perceptie si detectia vorbirii.

\end{document}
